\section{Ograniczenia pracy}

Mimo wyraźnych wyników badania, należy uwzględnić kilka ograniczeń przy ich interpretacji.

Eksperymenty przeprowadzono na trzech taliach referencyjnych. To ustawienie zapewnia porównywalność wyników i kontrolę warunków, ale nie obejmuje pełnej zmienności gry Hearthstone. Skuteczność uzyskanych agentów należy traktować jako potwierdzoną dla badanego przekroju talii, a nie dla całego środowiska gry.

Przy wyrównaniu budżetu obliczeniowego według czasu uruchomienia (sekcja~\ref{sec:budget}) liczba powtórzeń wariantu głębokościowego została świadomie ograniczona do jednego na schemat optymalizacji --- pojedyncze uruchomienie trwa około \(3{,}5\) razy dłużej niż pozostałe konfiguracje (sekcja~\ref{sec:setup}). To utrzymuje porównywalny nakład obliczeniowy między wariantami, ale ogranicza siłę statystyczną wniosków dotyczących tego wariantu względem konfiguracji uruchamianych trzykrotnie.

Faza III opiera się na jednym zewnętrznym punkcie odniesienia --- agencie Sebastiana Millera z edycji 2020. Taka walidacja ma wysoką wartość poznawczą, ale nie jest pełną replikacją rankingu konkursowego, który obejmuje szerszy zbiór stylów gry i implementacji.

Analiza wyników ma przede wszystkim charakter zbiorczy. Oceniano odsetki zwycięstw, przebiegi jakości populacji i profile wag, ale nie wykonano analizy decyzji na poziomie pojedynczych tur i sekwencji akcji. Część interpretacji mechanizmów działania agentów pozostaje więc na poziomie dobrze uzasadnionej hipotezy.

